% GENERATED FILE. Edit canonical lessons, then run npm run generate:books.
% canonical-source-hash: fnv1a64:19a6de19d0ca571a
% canonical-lessons: JA-C402-kankou, JA-C402-kankoukyaku, JA-C402-meisho, JA-C402-meibutsu, JA-C402-meguru

\chapter{A Day of Sightseeing}
\label{ch:ja-a-day-of-sightseeing}
\hlchaptermodality{402}

\begin{chapteropening}
\textbf{By the end of this chapter:} \emph{I can plan a day of sightseeing.}

Complete the last lesson of chapter 402: I can plan a day of sightseeing.
\end{chapteropening}

\section[\texorpdfstring{\ja{かんこう} (kankō)}{かんこう (kankō)}]{\ja{かんこう} (kankō) — tourism; sightseeing}

\label{lesson:JA-C402-kankou}



\begin{quote}
Before the new one: say the Japanese for to change one's clothes, then the Japanese for to entertain guests.
\end{quote}

\subsection*{You'll want to know: \ja{かんこう}}
\textbf{\ja{かんこう}} — \emph{kankō} — \textquotedblleft{}tourism; sightseeing\textquotedblright{}.

\begin{quote}
\textbf{\ja{この} \ja{しまに} \ja{かんこうに} \ja{いきました。}} — \emph{kono shima ni kankō ni ikimashita} — \textquotedblleft{}We went sightseeing on this island.\textquotedblright{}
\end{quote}

\begin{cousinweb}
\textbf{\ja{かんこう}} (\emph{kankō}) is Sino-Japanese (kango): two Chinese characters read with their on'yomi, readings borrowed from Chinese. They are \emph{kan}, \textquotedblleft{}view\textquotedblright{} (Middle Chinese \emph{kwan}) and \emph{kō}, \textquotedblleft{}light\textquotedblright{} (\emph{kwang}) (KANJIDIC2; Middle Chinese in Baxter's transcription).
\end{cousinweb}

\begin{grammarlens}[title={What you've built}]
One more word for getting around.
\end{grammarlens}

\subsection*{Your turn}
\begin{itemize}
\raggedright
  \item \emph{Say it:} \emph{kankō}
  \item \emph{Say it:} \emph{kankō}, once more
  \item \emph{Say it:} \emph{kono shima ni kankō ni ikimashita} — We went sightseeing on this island.
  \item \emph{Recall:} say the Japanese for to change one's clothes, then the Japanese for to entertain guests, then say \emph{kankō} again
\end{itemize}

\begin{tcolorbox}[breakable,colback=teal!4,colframe=teal!35!black,title={Before you move on}]
What does \ja{かんこう} mean? (\textquotedblleft{}Tourism; sightseeing\textquotedblright{}.) Say it once more.
\end{tcolorbox}
\section[\texorpdfstring{\ja{かんこうきゃく} (kankōkyaku)}{かんこうきゃく (kankōkyaku)}]{\ja{かんこうきゃく} (kankōkyaku) — a tourist}

\label{lesson:JA-C402-kankoukyaku}



\begin{quote}
Before the new one: say the Japanese for to entertain guests, then the Japanese for tourism.
\end{quote}

\subsection*{You'll want to know: \ja{かんこうきゃく}}
\textbf{\ja{かんこうきゃく}} — \emph{kankōkyaku} — \textquotedblleft{}a tourist\textquotedblright{}.

\begin{quote}
\textbf{\ja{なつは} \ja{かんこうきゃくが} \ja{おおいです。}} — \emph{natsu wa kankōkyaku ga ōi desu} — \textquotedblleft{}There are a lot of tourists in summer.\textquotedblright{}
\end{quote}

\begin{cousinweb}
\textbf{\ja{かんこうきゃく}} (\emph{kankōkyaku}) is Sino-Japanese (kango): three Chinese characters read with their on'yomi, readings borrowed from Chinese. They are \emph{kan}, \textquotedblleft{}view\textquotedblright{} (Middle Chinese \emph{kwan}), \emph{kō}, \textquotedblleft{}light\textquotedblright{} (\emph{kwang}) and \emph{kyaku}, \textquotedblleft{}guest\textquotedblright{} (\emph{khaek}) (KANJIDIC2; Middle Chinese in Baxter's transcription).
\end{cousinweb}

\begin{grammarlens}[title={What you've built}]
One more word for getting around.
\end{grammarlens}

\subsection*{Your turn}
\begin{itemize}
\raggedright
  \item \emph{Say it:} \emph{kankōkyaku}
  \item \emph{Say it:} \emph{kankōkyaku}, once more
  \item \emph{Say it:} \emph{natsu wa kankōkyaku ga ōi desu} — There are a lot of tourists in summer.
  \item \emph{Recall:} say the Japanese for to entertain guests, then the Japanese for tourism, then say \emph{kankōkyaku} again
\end{itemize}

\begin{tcolorbox}[breakable,colback=teal!4,colframe=teal!35!black,title={Before you move on}]
What does \ja{かんこうきゃく} mean? (\textquotedblleft{}A tourist\textquotedblright{}.) Say it once more.
\end{tcolorbox}
\section[\texorpdfstring{\ja{めいしょ} (meisho)}{めいしょ (meisho)}]{\ja{めいしょ} (meisho) — a famous sight; a place of interest}

\label{lesson:JA-C402-meisho}



\begin{quote}
Before the new one: say the Japanese for tourism, then the Japanese for a tourist.
\end{quote}

\subsection*{You'll want to know: \ja{めいしょ}}
\textbf{\ja{めいしょ}} — \emph{meisho} — \textquotedblleft{}a famous sight; a place of interest\textquotedblright{}.

\begin{quote}
\textbf{\ja{この} \ja{まちの} \ja{めいしょを} \ja{まわりました。}} — \emph{kono machi no meisho o mawarimashita} — \textquotedblleft{}We went round the sights of this town.\textquotedblright{}
\end{quote}

\begin{cousinweb}
\textbf{\ja{めいしょ}} (\emph{meisho}) is Sino-Japanese (kango): two Chinese characters read with their on'yomi, readings borrowed from Chinese. They are \emph{mei}, \textquotedblleft{}noted\textquotedblright{} (Middle Chinese \emph{mjieng}) and \emph{sho}, \textquotedblleft{}place\textquotedblright{} (\emph{srjoX}) (KANJIDIC2; Middle Chinese in Baxter's transcription).
\end{cousinweb}

\begin{grammarlens}[title={What you've built}]
One more word for getting around.
\end{grammarlens}

\subsection*{Your turn}
\begin{itemize}
\raggedright
  \item \emph{Say it:} \emph{meisho}
  \item \emph{Say it:} \emph{meisho}, once more
  \item \emph{Say it:} \emph{kono machi no meisho o mawarimashita} — We went round the sights of this town.
  \item \emph{Recall:} say the Japanese for tourism, then the Japanese for a tourist, then say \emph{meisho} again
\end{itemize}

\begin{tcolorbox}[breakable,colback=teal!4,colframe=teal!35!black,title={Before you move on}]
What does \ja{めいしょ} mean? (\textquotedblleft{}A famous sight; a place of interest\textquotedblright{}.) Say it once more.
\end{tcolorbox}
\section[\texorpdfstring{\ja{めいぶつ} (meibutsu)}{めいぶつ (meibutsu)}]{\ja{めいぶつ} (meibutsu) — a local speciality}

\label{lesson:JA-C402-meibutsu}



\begin{quote}
Before the new one: say the Japanese for a tourist, then the Japanese for a famous sight.
\end{quote}

\subsection*{You'll want to know: \ja{めいぶつ}}
\textbf{\ja{めいぶつ}} — \emph{meibutsu} — \textquotedblleft{}a local speciality\textquotedblright{}.

\begin{quote}
\textbf{\ja{この} \ja{まちの} \ja{めいぶつは} \ja{さかなです。}} — \emph{kono machi no meibutsu wa sakana desu} — \textquotedblleft{}Fish is this town's speciality.\textquotedblright{}
\end{quote}

\begin{cousinweb}
\textbf{\ja{めいぶつ}} (\emph{meibutsu}) is Sino-Japanese (kango): two Chinese characters read with their on'yomi, readings borrowed from Chinese. They are \emph{mei}, \textquotedblleft{}noted\textquotedblright{} (Middle Chinese \emph{mjieng}) and \emph{butsu}, \textquotedblleft{}thing\textquotedblright{} (\emph{mjut}) (KANJIDIC2; Middle Chinese in Baxter's transcription).
\end{cousinweb}

\begin{grammarlens}[title={What you've built}]
One more word for getting around.
\end{grammarlens}

\subsection*{Your turn}
\begin{itemize}
\raggedright
  \item \emph{Say it:} \emph{meibutsu}
  \item \emph{Say it:} \emph{meibutsu}, once more
  \item \emph{Say it:} \emph{kono machi no meibutsu wa sakana desu} — Fish is this town's speciality.
  \item \emph{Recall:} say the Japanese for a tourist, then the Japanese for a famous sight, then say \emph{meibutsu} again
\end{itemize}

\begin{tcolorbox}[breakable,colback=teal!4,colframe=teal!35!black,title={Before you move on}]
What does \ja{めいぶつ} mean? (\textquotedblleft{}A local speciality\textquotedblright{}.) Say it once more.
\end{tcolorbox}
\section[\texorpdfstring{\ja{めぐる} (meguru)}{めぐる (meguru)}]{\ja{めぐる} (meguru) — to tour; to go around}

\label{lesson:JA-C402-meguru}



\begin{quote}
Before the new one: say the Japanese for a famous sight, then the Japanese for a local speciality.
\end{quote}

\subsection*{You'll want to know: \ja{めぐる}}
\textbf{\ja{めぐる}} — \emph{meguru} — \textquotedblleft{}to tour; to go around\textquotedblright{}.

\begin{quote}
\textbf{\ja{じてんしゃで} \ja{しまの} \ja{めいしょを} \ja{めぐりました。}} — \emph{jitensha de shima no meisho o megurimashita} — \textquotedblleft{}We toured the island's sights by bicycle.\textquotedblright{}
\end{quote}

Polite \textbf{\ja{めぐります}}, \emph{megurimasu}; a godan verb. Past (polite): \textbf{\ja{かんこうきゃくが} \ja{まちを} \ja{めぐりました。}}, \emph{kankōkyaku ga machi o megurimashita}.

\begin{cousinweb}
\textbf{\ja{めぐる}} (\emph{meguru}) is a native Japanese word (wago), inherited rather than borrowed from Chinese; its kanji was attached later as a kun'yomi (JMdict).
\end{cousinweb}

\begin{grammarlens}[title={What you've built}]
One more word for getting around.
\end{grammarlens}

\subsection*{Your turn}
\begin{itemize}
\raggedright
  \item \emph{Say it:} \emph{meguru}
  \item \emph{Say it:} \emph{meguru}, once more
  \item \emph{Say it:} \emph{jitensha de shima no meisho o megurimashita} — We toured the island's sights by bicycle.
  \item \emph{Recall:} say the Japanese for a famous sight, then the Japanese for a local speciality, then say \emph{meguru} again
\end{itemize}

\begin{tcolorbox}[breakable,colback=teal!4,colframe=teal!35!black,title={Before you move on}]
What does \ja{めぐる} mean? (\textquotedblleft{}To tour; to go around\textquotedblright{}.) Say it once more.
\end{tcolorbox}
